% Marcel Harrer — CV (international / senior AI engineer, English)
% Style: classic LaTeX look (Computer Modern), single column, ATS-safe.
\documentclass[10pt,a4paper]{article}

\usepackage[a4paper,top=0.9cm,bottom=0.9cm,left=1.5cm,right=1.5cm]{geometry}
\usepackage[T1]{fontenc}
\usepackage[utf8]{inputenc}
\usepackage{enumitem}
\usepackage{titlesec}
\usepackage[hidelinks]{hyperref}
\usepackage{amsmath}

\pagestyle{empty}
\setlength{\parindent}{0pt}

% --- Section style: bold heading with a rule underneath (matches original) ---
\titleformat{\section}{\large\bfseries}{}{0pt}{}[\vspace{1pt}\titlerule]
\titlespacing*{\section}{0pt}{5pt}{2.5pt}

% --- Helpers ---
\newcommand{\sep}{\,$\cdot$\,}
\newcommand{\role}[4]{% {Title}{Company}{Location}{Dates}
  \vspace{2pt}%
  {\textbf{#1}, #2 \hfill \textit{#3\ifx&#3&\else\sep\fi#4}}\par}
\newlist{cvitems}{itemize}{1}
\setlist[cvitems]{label=\textbullet, leftmargin=14pt, itemsep=0.5pt, parsep=0pt, topsep=1.5pt}

\begin{document}

% ================= HEADER =================
\begin{center}
  {\LARGE\bfseries Marcel Harrer}\\[3pt]
  AI Engineer \sep LLM Systems, RAG \& Multi-Agent Architectures\\
  Vienna, Austria \sep open to remote \sep +43\,660\,630\,2084 \sep
  \href{mailto:harrer.marcel@gmail.com}{harrer.marcel@gmail.com}\\
  \href{https://github.com/MarcelHarrer}{github.com/MarcelHarrer} \sep
  \href{https://linkedin.com/in/marcel-harrer}{linkedin.com/in/marcel-harrer}
\end{center}

% ================= SUMMARY =================
\section{Summary}
AI Engineer specializing in production LLM systems, RAG, and multi-agent
architectures, with further depth in computer vision and NLP. Currently
architecting agent systems that autonomously process $\sim$10 million
government emails a year across 50 departments. M.Sc.\ in Applied Mathematics
--- pairs hands-on engineering with mathematical rigor, incl.\ research in
preparation for publication.

% ================= EXPERIENCE =================
\section{Experience}

\role{AI Engineer}{Fivesquare}{Remote}{Apr 2026 -- Present}
\begin{cvitems}
  \item Shipped a multi-agent system that automatically handles $\sim$10 million
        incoming government emails a year across 50 departments, replacing fully
        manual processing.
  \item Architected a shared orchestrator routing to a dedicated agent per
        department, each learning new skills from production experience (agentic
        context engineering) instead of requiring constant retraining.
  \item Own production LLM and agent systems in Python end to end --- design,
        deployment, and operation --- on the core AI team.
\end{cvitems}

\role{AI Engineer (Freelance)}{Independent}{Remote}{Mar 2025 -- Dec 2025}
\begin{cvitems}
  \item Delivered a Text-to-SQL agent (LangGraph) that lets non-technical staff
        query PostgreSQL databases in plain English, removing their dependence on
        engineers for routine reporting.
  \item Designed the agent's modular tool interfaces --- database querying,
        reasoning over structured data, and email workflows --- so it completes
        multi-step requests end to end rather than one-off queries.
  \item Tuned orchestration and prompt strategies for reliable, repeatable
        production results --- not just a working demo.
\end{cvitems}

\role{AI Research Engineer}{ImageTwin}{}{Jan 2024 -- Dec 2025}
{\footnotesize\itshape Part-time in 2025, concurrent with the freelance role above.}
\begin{cvitems}
  \item Built computer-vision models that flag manipulated and duplicated images
        in scientific figures by matching them against a 100-million-image
        dataset, catching integrity issues nearly impossible to spot by eye.
  \item Implemented keypoint-based image matching (OpenCV SIFT/SURF) and trained
        LightGlue models to reliably detect reused or altered image regions.
  \item Co-built an end-to-end pipeline to detect AI-generated images --- from
        labeling training data through automated evaluation --- as newer fakes
        began slipping past existing checks.
\end{cvitems}

\role{AI Engineer}{Financial Market Authority (FMA)}{Vienna, Austria}{Dec 2022 -- Dec 2023}
\begin{cvitems}
  \item Built a RAG search assistant (FAISS + transformer embeddings) giving
        regulators instant, source-backed answers from thousands of
        German-language regulatory documents.
  \item Made large regulatory archives searchable by meaning rather than exact
        wording (SBERT similarity), surfacing relevant rules staff would otherwise
        miss and clustering related texts to speed up review.
  \item Built dashboards in R that turned complex Austrian insurance-sector data
        into clear visuals for non-technical decision-makers.
\end{cvitems}

% ================= RESEARCH =================
\section{Research \& Selected Projects}

\textbf{B-Spline Networks and a Universal Approximation Theorem}
\textit{(in preparation, with Prof.\ Arnold Neumaier)}\\
Introduced ``spline networks'' (single-hidden-layer networks with adaptive cubic
splines over random ridge directions) and proved a quantitative universal
approximation theorem bounding the expected $L^2$ error by $O(1/m + h^8)$;
training reduces to sparse linear least squares, solved with conjugate gradients in Python.

\vspace{2pt}
\textbf{Master's Thesis: Accelerating Gradient Descent via L-Smoothness}\\
Adaptive optimizer that leverages L-smoothness to dynamically tune step sizes,
with provable convergence guarantees and faster convergence than fixed-step methods.

% ================= EDUCATION =================
\section{Education}
{\textbf{M.Sc.\ in Applied Mathematics}, University of Vienna \hfill \textit{Oct 2020 -- Nov 2023}}\\
{\textbf{B.Sc.\ in Mathematics}, University of Vienna \hfill \textit{Oct 2016 -- Sep 2020}}

% ================= SKILLS =================
\section{Technical Skills}
\textbf{Languages:} Python, SQL\\
\textbf{LLMs \& Agents:} RAG, multi-agent orchestration, LangGraph, LangChain,
MCP (Model Context Protocol), prompt engineering, Hugging Face Transformers,
vLLM, Ollama, AWS Bedrock\\
\textbf{ML \& Computer Vision:} PyTorch, TensorFlow, OpenCV (SIFT/SURF),
LightGlue, time-series forecasting\\
\textbf{Retrieval \& Data:} FAISS, Qdrant, pgvector, semantic search, PostgreSQL\\
\textbf{MLOps \& Infrastructure:} FastAPI, Docker, Kubernetes, CI/CD,
AWS (SageMaker, Lambda, S3, EC2), Weights \& Biases, Langfuse\\
\textbf{Spoken Languages:} German (native), English (bilingual, C2),
Spanish (advanced, C1)

\end{document}
